\noindent\textbf{ON NETWORK ADJUSTMENT.}\hfill319
\begin{center}\textbf{[10.]}\par\medskip
\textbf{[Correspondence on Network Adjustment.]}\end{center}

\begin{center}\textsc{Gauss} to \textsc{Olbers}.\quad Göttingen, 2 November 1823.\end{center}

.\ .\ .\ .\ .\ .\ I have now completed the laborious adjustment of all my
measurements from 1821---1823, as far as the principal triangles are concerned,
so that both the sums of the angles in the individual triangles and the ratios of
the sides in the crossed quadrilaterals and pentagons now agree exactly, in
accordance with the strict principles of the calculus of probabilities, sine
ira et studio, and without any arbitrariness. It does not seem that the latter
can be said of the measurements of other geodesists. In the entire system
there are 76 directions, i.e. 38, each measured both ways. For \textls[120]{none}
of them did the adjustment amount to $1^{\prime\prime}$. The largest is
$0^{\prime\prime}{,}813$ for the direction from Nindorf to Hamburg, where
pointing at the Michaelisthurm was always very difficult in the smoke-like haze
and with the phase error. The next largest is $0^{\prime\prime}{,}788$ for the
direction from Lüneburg to Wilsede. There the target was a heliotrope light,
but the mounting on a wooden tripod in the lantern of the tower was certainly
less solid than on stone pedestals, and the observer's weight may have affected
the instrument as he changed position. The \textls[120]{mean} error of all the
directions, in the sense used in my Theoria Combinationis, is $0^{\prime\prime}{,}486$.

In the entire system there are 26 triangles in which I measured all
the angles. The greatest error in the sum of the angles is now $2^{\prime\prime}{,}175$
in the triangle Nindorf-Hamburg-Timpenberg, for which the direction from
Nindorf to Hamburg may be chiefly to blame. The next greatest error, as
mentioned above, is in the triangle Brocken-Hohehagen-Hils. It amounts
to $1^{\prime\prime}{,}806$, and the direction from Hils to Brocken is now
likely to be chiefly to blame.
\par\medskip
.\ .\ .\ .\ .\ .

\clearpage
\noindent320\hfill\textbf{CORRESPONDENCE.}
\begin{center}\textsc{Gauss} to \textsc{Olbers}.\quad Zeven, 6 July 1824.\end{center}

\begingroup
\makeatletter
\renewcommand{\@makefnmark}{\hbox{\normalfont\textsuperscript{*})}}
\makeatother
.\ .\ .\ .\ .\ .\ The new measurements from 1823 [have] produced some
changes in the entire system, albeit very small ones\footnote{In my current calculation I have also used \textsc{Walbeck}'s flattening \textls[120]{exactly}, whereas in the previous one, as you know, an oversight led to a discrepancy, albeit an extremely small one.}.
\par\medskip
.\ .\ .\ .\ .\ .
\endgroup
